\chapter{IMPLEMENTATION}
\thispagestyle{empty}

This chapter covers the technology selection, the repository structure, and
the implementation of the backend, frontend, and infrastructure layers of
the platform.

\clearpage

\section{Technology Selection}

The stack was selected against the constraints of a multi-tenant SaaS
platform with document processing and AI generation. Table~\ref{tab:tech}
lists the technologies with the reasons for their selection; references are
provided in the bibliography.

\begin{table}[H]
\centering
\caption{Technology selection.}
\label{tab:tech}
\small
\begin{tabular}{@{}p{4.6cm} p{\dimexpr\textwidth-4.6cm-12pt\relax}@{}}
\toprule
\textbf{Technology} & \textbf{Role and rationale} \\
\midrule
NestJS / TypeScript \cite{nestjs} & Modular monolith API with typed, linted
codebase and explicit module boundaries. \\
Next.js \cite{nextjs} / Node.js \cite{nodejs} & Server-rendered web application
sharing the TypeScript contract packages. \\
PostgreSQL \cite{postgresql} + Drizzle ORM \cite{drizzle} & Single primary
database with JSONB for flexible payloads; typed migrations. \\
RabbitMQ \cite{rabbitmq} & Asynchronous task distribution with durable queues
for material and AI work. \\
Python 3.12 \cite{python} + FastAPI \cite{fastapi} & OCR service and worker
processes with type-checked, linted code. \\
Pydantic \cite{pydantic} & Worker-boundary validation of job payloads and AI
output. \\
PyMuPDF \cite{pymupdf} + PaddleOCR \cite{ppocr,ppocrv3} & Embedded PDF text
extraction with per-page OCR fallback. \\
Zod & API/contract validation shared between web and API. \\
pnpm workspaces + Turborepo & Monorepo build and type-check orchestration. \\
\bottomrule
\end{tabular}
\end{table}

\section{Repository Structure}

The monorepo groups the deployable applications and the shared packages:

\begin{lstlisting}[basicstyle=\ttfamily\footnotesize]
apps/
  api/        # NestJS modular monolith (global /api/v1 prefix)
  web/        # Next.js 15 application
  ocr/        # FastAPI OCR service
  workers/    # Python workers (material, AI) + standalone ocr-worker
packages/
  contracts/  # Zod DTO and response schemas (shared wire contract)
  database/   # Drizzle schema, migrations, seed
  auth/       # authentication/authorization utilities
  ai/         # AI integration types and utilities
  shared/     # shared utilities (e.g. email normalization)
infrastructure/
  compose/    # Dockerfiles, requirements lock, compose files
  nginx/      # internal reverse proxy configuration
docs/         # API reference, architecture, status, tasks, diagrams
\end{lstlisting}

\section{Backend Implementation}

\subsection{Identity, tenancy, and authorization}

Authentication is cookie-based. Login and refresh set httpOnly
\texttt{access\_\hspace{0pt}token} and \texttt{refresh\_\hspace{0pt}token}
cookies; the refresh token
is hashed into \texttt{auth\_\hspace{0pt}sessions} and rotated on every
refresh with a
race-handling guard. A CSRF double-submit token is issued alongside the
session. Every guarded route runs the access-token guard, the tenant guard
(validating the \texttt{x-institute-id} header against an active membership),
and the roles guard. Queries are consistently tenant-scoped in application
code; for example the jobs service filters every lookup by both job id and
institute:

\begin{lstlisting}[basicstyle=\ttfamily\footnotesize]
const job = await db.query.jobs.findFirst({
  where: and(eq(jobs.id, jobId), eq(jobs.instituteId, instituteId)),
});
\end{lstlisting}

\subsection{Asynchronous jobs}

Long-running work is never executed in the API process. The API inserts a row
into the \texttt{jobs} table and publishes a message to RabbitMQ; the worker
that consumes the message runs the work and reports the terminal state back
to the API, which updates the job row. Job lifecycle is
\texttt{queued} $\rightarrow$ \texttt{processing} $\rightarrow$
\texttt{completed} \texttt{|} \texttt{failed}, with a racing-safe cancel path
(\texttt{cancelling} $\rightarrow$ \texttt{cancelled}). Active AI-generation
jobs are de-duplicated by operation and source.

\subsection{Materials and the OCR pipeline}

Materials accept text input or uploaded files (PDF, images); large uploads
are transferred in chunks over an XMLHttpRequest-based client to avoid
tunnel origin timeouts. The OCR pipeline extracts embedded text with PyMuPDF
first and falls back to PaddleOCR per page; normalization is deterministic
and intentionally never invents text:

\begin{lstlisting}[language=Python,basicstyle=\ttfamily\footnotesize]
text = text.replace("\x0c", "\n").replace("\xa0", " ")
text = text.replace("\r\n", "\n").replace("\r", "\n")
text = "".join(ch for ch in text
               if ch == "\n" or ch == "\t"
               or unicodedata.category(ch)[0] != "C")
lines = [re.sub(r"[ \t]+", " ", line).strip()
         for line in text.split("\n")]
\end{lstlisting}

Processing is distributed. The API hosts an OCR coordinator that materializes
page-range chunks, leases them to workers, retries failed or timed-out chunks,
and settles cancelled jobs; the workers may run inside the stack or as
standalone devices that connect to the API over HTTPS. Per-page correction
records are stored for inspection.

\subsection{Material Intelligence and extraction}

A deterministic enhancer cleans and structures the extracted text, then
segments it into logical regions (heading, paragraph, table, figure and
similar kinds). Each segment is mapped to syllabus entities (subject, chapter,
topic, or unit) with a confidence value; mappings carry their full ancestor
context and are governed by a type-aware check constraint. Enhancement
results are append-only and versioned, and a source hash keeps re-runs
idempotent.

Deterministic extractors recover paper patterns and question candidates from
text, PDF, and image sources, and always produce explicit ambiguity reports
that a human resolves instead of guessing silently.

\subsection{AI generation}

All AI traffic flows through an internal OpenAI-compatible gateway
(OmniRoute); no local model and no direct cloud SDK calls are used. The AI
worker operates on a dedicated RabbitMQ queue with bounded concurrency.
Long sources are split deterministically into overlapping chunks; each chunk
is sent to the provider, its output is validated against Pydantic schemas,
invalid chunks are re-requested after a transient-aware retry, and results
are aggregated deterministically. Generation operations cover notes,
summaries, flashcards, important concepts, complete content packages
(including Cornell notes), question-bank batches, paper-pattern blueprints,
starter material, and syllabus analysis.

\subsection{Examinations, attempts, and practice}

Examinations assemble question papers against approved paper patterns with
coverage gating. Attempts begin by snapshotting the question set into an
immutable form, enforce start/end deadlines server-side, validate answers per
question type, save idempotently, and evaluate objective question types
(MCQ, true/false, fill-in-the-blank, numerical, matching) deterministically
server-side. Result analytics summarize performance. Practice is deliberately
ungraded: per-item instant feedback with reveal, and no score.

\subsection{Export}

Content, question banks, assessments with results, paper patterns, and
question papers are rendered to DOCX (via document templates), PDF (via
server-side Chromium through Puppeteer), and XLSX spreadsheets, each with a
preview path in the web application.

\section{Frontend Implementation}

The web application is built with Next.js 15 using strict TypeScript, Tailwind
CSS, and a shared design system of shadcn/ui components layered on Radix
primitives. A single typed API client adds cookie credentials, the tenant
header, CSRF headers, automatic job polling (\texttt{waitForJob}), and chunked
uploads. Role-aware navigation presents teacher and student workspaces, and
every asynchronous surface renders loading, success, empty, error, and retry
states. The web app shares the Zod contracts with the API through the
contracts package, so payload shapes cannot drift between layers.

\section{Notes on the Approved Proposal}

The implementation follows the approved proposal's scope and requirements.
Two points are recorded here for clarity:

\begin{enumerate}
  \item \textbf{Async workers.} As proposed, asynchronous work is distributed
        over RabbitMQ; the Python workers consume the queues directly over
        AMQP with the \texttt{pika} client instead of a task-queue framework.
        The messaging design is unchanged.
  \item \textbf{Caching layer.} Redis is declared in the compose stack, but no
        application code currently uses it; it is therefore recorded as a
        limitation rather than claimed as implemented.
\end{enumerate}

Flashcards and notes are implemented as content types of the content module
with type-specific payload contracts, which matches the proposal's intent for
study content.
% ---------- API Surface (folded from former appendix-api) ----------
\section{API Surface}

All routes are served under the global \texttt{/api/v1} prefix. The table
lists representative endpoints per module; it is not an exhaustive
enumeration: representative behaviour is listed per module, served under the
global \texttt{/api/v1} prefix.

\begin{table}[H]
\centering
\caption{Representative API endpoints.}
\label{tab:api}
\small
\begin{tabular}{@{}p{3.2cm}p{5.2cm}p{\dimexpr\textwidth-8.4cm-24pt\relax}@{}}
\toprule
\textbf{Module} & \textbf{Endpoint} & \textbf{Behaviour} \\
\midrule
identity & POST /auth/login & Authenticate, set session cookies \\
identity & POST /auth/refresh & Rotate refresh token, new access token \\
identity & GET /auth/me & Current user and membership context \\
tenancy & GET /memberships & Institute picker for the user \\
users & GET, POST /users & List and create institute users \\
users & PATCH /users/:id/status & Activate or deactivate a member \\
jobs & POST /jobs & Create and publish an async job \\
jobs & POST /jobs/:id/retry, /cancel & Retry or cancel a job \\
academic & GET/POST/... /academic/subjects & CRUD on subjects \\
academic & .../subjects/:id/chapters & CRUD on chapters and topics \\
syllabus & POST /academic/subjects/:id/syllabus/... & Generate, confirm, lock syllabi \\
materials & POST /materials & Create material (text) \\
materials & POST /materials (upload) & Upload a file, chunked \\
materials & POST /materials/:id/process, /retry & Process or retry OCR \\
materials & GET /materials/:id/enhancement & Material Intelligence result \\
questions & GET/POST /questions & Question bank CRUD \\
questions & POST /questions/generate & AI question generation \\
question-types & GET/POST /question-types & Data-driven question types \\
paper-patterns & POST /paper-patterns/generate & AI blueprint generation \\
paper-patterns & POST /paper-patterns/analyze & Analyse a pattern source \\
question-papers & POST /question-papers & Create a question paper \\
question-papers & POST /question-papers/:id/select-from-pattern & Auto-select questions \\
examinations & POST /examinations (assessments) & Create an assessment \\
examinations & POST /assessments/:id/publish & Publish an assessment \\
attempts & POST /attempts/start & Snapshot and start an attempt \\
attempts & POST /attempts/:id/submit & Submit and evaluate \\
attempts & GET /attempts/:id/result, /analytics & Results and analytics \\
practice & POST /practice/sessions & Start a practice session \\
export & GET/POST /export/\ldots & DOCX, PDF, XLSX export \\
health & GET /health & Readiness probe \\
\bottomrule
\end{tabular}
\end{table}
